%% ============================================================
%% AULA 13 — Aplicações: classificação em escala e projeto de
%%           redes
%% GBC073 — Inteligência Computacional (FACOM/UFU)
%%
%% Esqueleto com esboço de conteúdo — a desenvolver.
%% Ementa (ficha SEI 5116656): item 2 — aplicações: classificação;
%% projeto de redes.
%%
%% Compilar:  latexmk -xelatex aula13.tex   (duas passadas!)
%% ============================================================
\documentclass[aspectratio=169, 11pt]{beamer}

\makeatletter
\def\input@path{{./}{../}}
\makeatother

\usetheme{InteligenciaComputacional}   % ANTES do babel: fontspec vem primeiro
\usepackage{babel}
\babelprovide[import, main]{brazilian}

\title[Aula 13 — Classificação em escala]{Aula 13: Aplicações — classificação em escala e projeto de redes}
\subtitle[GBC073 — Inteligência Computacional]{Visão \textbullet\ áudio \textbullet\ saúde \textbullet\ boas práticas}
\author[Prof.\ Marcelo Albertini]{Prof.\ Marcelo Keese Albertini}
\institute{GBC073 — Inteligência Computacional \textbullet\ Universidade Federal de Uberlândia}
\date{2026}

\begin{document}

% ------------------------------------------------------------ capa
\begin{frame}[plain, noframenumbering]
  \titlepage
\end{frame}

% ------------------------------------------------------------ roteiro
\begin{frame}{Roteiro da aula}
  \begin{itemize}\setlength{\itemsep}{0.5ex}
    \item<1-> \textbf{Classificação em escala (visão computacional)}
    \item<2-> \textbf{Áudio, saúde e texto}
    \item<3-> \textbf{Projeto de redes: do problema à arquitetura}
    \item<4-> \textbf{Boas práticas: baselines, ablações, logs}
  \end{itemize}
\end{frame}

% ------------------------------------------------------------
\section{Classificação em escala}

\begin{frame}{Classificação em escala (visão computacional)}
  \begin{itemize}\setlength{\itemsep}{0.4ex}
    \item De 10 classes (CIFAR-10) a milhares (ImageNet) e milhões (retrieval).
    \item A jornada do problema: dados, arquitetura, treino, avaliação.
    \item Métricas que importam: acurácia não conta a história inteira.
    \item Detecção e segmentação como parentes próximos da classificação.
  \end{itemize}
\end{frame}

% ------------------------------------------------------------
\section{Áudio, saúde e texto}

\begin{frame}{Áudio, saúde e texto}
  \begin{itemize}\setlength{\itemsep}{0.4ex}
    \item Áudio: espectrogramas + CNNs; voz, música, sons do ambiente.
    \item Saúde: imagens médicas, sinais, prontuários — e o peso de errar.
    \item Texto: classificação de documentos, sentimento, moderação.
    \item O mesmo arcabouço (aulas 7--9), domínios diferentes — a
          representação é que muda.
  \end{itemize}
\end{frame}

% ------------------------------------------------------------
\section{Projeto de redes}

\begin{frame}{Projeto de redes: do problema à arquitetura}
  \begin{itemize}\setlength{\itemsep}{0.4ex}
    \item Topologia: quantas camadas, quantos neurônios, que ativações.
    \item Parâmetros: taxa de aprendizado, tamanho de lote, épocas.
    \item Modos de treinamento: do zero, transferência, ajuste fino.
    \item A ficha pede ``projeto de redes'' — a resposta moderna é
          \emph{começar de algo que já funciona}.
  \end{itemize}
\end{frame}

% ------------------------------------------------------------
\section{Boas práticas}

\begin{frame}{Boas práticas: baselines, ablações, logs}
  \begin{itemize}\setlength{\itemsep}{0.4ex}
    \item \emph{Baseline} honesto primeiro: o método mais simples que já
          resolve o problema.
    \item Ablações: desligar uma peça por vez para medir a contribuição de
          cada uma.
    \item Registrar tudo: semente, versão do código, hiperparâmetros.
    \item Só melhore o que a métrica mede — e desconfie da métrica.
  \end{itemize}
  \begin{ideiabox}
    Um experimento reprodutível vale mais que um resultado impressionante:
    quem não registra as condições, não consegue nem repetir o próprio
    achado.
  \end{ideiabox}
\end{frame}

\end{document}
